% GENERATED FILE. Edit canonical lessons, then run npm run generate:books.
% canonical-source-hash: fnv1a64:a777d150f76626ce
% canonical-lessons: FR-C456-ruisseau, FR-C456-etang, FR-C456-marais, FR-C456-patauger, FR-C456-ruisseler

\chapter{Ruisseaux et marais}
\label{ch:fr-ruisseaux-et-marais}
\hlchaptermodality{456}

\begin{chapteropening}
\textbf{By the end of this chapter:} \emph{I can describe wet places: the stream, the pond, the marsh, splashing about and water running down.}

Complete the last lesson of chapter 456: I can describe wet places: the stream, the pond, the marsh, splashing about and water running down.
\end{chapteropening}

\section[le ruisseau]{le ruisseau — the stream, the brook}

\label{lesson:FR-C456-ruisseau}



\begin{quote}
Before the new one: say the French for to climb, then the French for to stumble.
\end{quote}

\subsection*{You'll want to know: le ruisseau}
\textbf{le ruisseau} — \textquotedblleft{}the stream, the brook\textquotedblright{}.

\begin{quote}
Un petit ruisseau traverse le jardin.
\end{quote}

\begin{quote}
\emph{A little stream runs across the garden.}
\end{quote}

Plural \emph{les ruisseaux}.

\begin{cousinweb}
\emph{ruisseau} goes back to Vulgar Latin \emph{rīvuscellus} (EtymDB), a small \emph{rīvus} 'stream'.
\end{cousinweb}

\begin{grammarlens}[title={What you've built}]
One more word for telling, explaining and describing at B1.
\end{grammarlens}

\subsection*{Your turn}
\begin{itemize}
\raggedright
  \item \emph{Say it:} \emph{le ruisseau}
  \item \emph{Say it:} \emph{le ruisseau}, once more
  \item \emph{Say it:} a sentence of your own with \emph{le ruisseau}
  \item \emph{Recall:} say the French for to climb, then the French for to stumble, then say \emph{le ruisseau} again
\end{itemize}

\begin{tcolorbox}[breakable,colback=teal!4,colframe=teal!35!black,title={Before you move on}]
What does le ruisseau mean? (\textquotedblleft{}The stream, the brook\textquotedblright{}.) Say it once more.
\end{tcolorbox}
\section[l'étang]{l'étang — the pond}

\label{lesson:FR-C456-etang}



\begin{quote}
Before the new one: say the French for to stumble, then the French for the stream.
\end{quote}

\subsection*{You'll want to know: l'étang}
\textbf{l'étang} — \textquotedblleft{}the pond\textquotedblright{}.

\begin{quote}
Des canards nagent sur l'étang.
\end{quote}

\begin{quote}
\emph{Some ducks are swimming on the pond.}
\end{quote}

A masculine noun. Bigger than a \emph{mare}, smaller than a \emph{lac}.

\begin{cousinweb}
\emph{étang} comes from Old French \emph{estanc} (EtymDB lists it), built on \emph{estanchier} 'to stop the flow of water', which is usually traced to a Vulgar Latin verb; the details are not settled.
\end{cousinweb}

\begin{grammarlens}[title={What you've built}]
One more word for telling, explaining and describing at B1.
\end{grammarlens}

\subsection*{Your turn}
\begin{itemize}
\raggedright
  \item \emph{Say it:} \emph{l'étang}
  \item \emph{Say it:} \emph{l'étang}, once more
  \item \emph{Say it:} a sentence of your own with \emph{l'étang}
  \item \emph{Recall:} say the French for to stumble, then the French for the stream, then say \emph{l'étang} again
\end{itemize}

\begin{tcolorbox}[breakable,colback=teal!4,colframe=teal!35!black,title={Before you move on}]
What does l'étang mean? (\textquotedblleft{}The pond\textquotedblright{}.) Say it once more.
\end{tcolorbox}
\section[le marais]{le marais — the marsh, the swamp}

\label{lesson:FR-C456-marais}



\begin{quote}
Before the new one: say the French for the stream, then the French for the pond.
\end{quote}

\subsection*{You'll want to know: le marais}
\textbf{le marais} — \textquotedblleft{}the marsh, the swamp\textquotedblright{}.

\begin{quote}
Beaucoup d'oiseaux vivent dans le marais.
\end{quote}

\begin{quote}
\emph{Many birds live in the marsh.}
\end{quote}

\emph{Les marais salants} — salt marshes.

\begin{cousinweb}
\emph{marais} comes from Old French \emph{mareis}, from Frankish \emph{marisk} 'marsh' (EtymDB).
\end{cousinweb}

\begin{grammarlens}[title={What you've built}]
One more word for telling, explaining and describing at B1.
\end{grammarlens}

\subsection*{Your turn}
\begin{itemize}
\raggedright
  \item \emph{Say it:} \emph{le marais}
  \item \emph{Say it:} \emph{le marais}, once more
  \item \emph{Say it:} a sentence of your own with \emph{le marais}
  \item \emph{Recall:} say the French for the stream, then the French for the pond, then say \emph{le marais} again
\end{itemize}

\begin{tcolorbox}[breakable,colback=teal!4,colframe=teal!35!black,title={Before you move on}]
What does le marais mean? (\textquotedblleft{}The marsh, the swamp\textquotedblright{}.) Say it once more.
\end{tcolorbox}
\section[patauger]{patauger — to splash about, to wade; to flounder}

\label{lesson:FR-C456-patauger}



\begin{quote}
Before the new one: say the French for the pond, then the French for the marsh.
\end{quote}

\subsection*{You'll want to know: patauger}
\textbf{patauger} — \textquotedblleft{}to splash about, to wade; to flounder\textquotedblright{}.

\begin{quote}
Les enfants pataugent dans le ruisseau.
\end{quote}

\begin{quote}
\emph{The children are splashing about in the stream.}
\end{quote}

Figurative: \emph{je patauge} — I am getting nowhere.

\begin{cousinweb}
\emph{patauger} is usually linked to \emph{patte} 'paw, foot', but the origin is not settled. EtymDB lists it without a source.
\end{cousinweb}

\begin{grammarlens}[title={What you've built}]
One more word for telling, explaining and describing at B1.
\end{grammarlens}

\subsection*{Your turn}
\begin{itemize}
\raggedright
  \item \emph{Say it:} \emph{patauger}
  \item \emph{Say it:} \emph{patauger}, once more
  \item \emph{Say it:} a sentence of your own with \emph{patauger}
  \item \emph{Recall:} say the French for the pond, then the French for the marsh, then say \emph{patauger} again
\end{itemize}

\begin{tcolorbox}[breakable,colback=teal!4,colframe=teal!35!black,title={Before you move on}]
What does patauger mean? (\textquotedblleft{}To splash about, to wade; to flounder\textquotedblright{}.) Say it once more.
\end{tcolorbox}
\section[ruisseler]{ruisseler — to stream, to run down (of water)}

\label{lesson:FR-C456-ruisseler}



\begin{quote}
Before the new one: say the French for the marsh, then the French for to splash about.
\end{quote}

\subsection*{You'll want to know: ruisseler}
\textbf{ruisseler} — \textquotedblleft{}to stream, to run down (of water)\textquotedblright{}.

\begin{quote}
La pluie ruisselle sur les vitres.
\end{quote}

\begin{quote}
\emph{The rain is streaming down the windows.}
\end{quote}

\emph{Le ruissellement} — run-off water.

\begin{cousinweb}
\emph{ruisseler} is built on \emph{ruissel}, an old form of \emph{ruisseau} 'stream' (EtymDB lists it), from Latin \emph{rīvus}.
\end{cousinweb}

\begin{grammarlens}[title={What you've built}]
One more word for telling, explaining and describing at B1.
\end{grammarlens}

\subsection*{Your turn}
\begin{itemize}
\raggedright
  \item \emph{Say it:} \emph{ruisseler}
  \item \emph{Say it:} \emph{ruisseler}, once more
  \item \emph{Say it:} a sentence of your own with \emph{ruisseler}
  \item \emph{Recall:} say the French for the marsh, then the French for to splash about, then say \emph{ruisseler} again
\end{itemize}

\begin{tcolorbox}[breakable,colback=teal!4,colframe=teal!35!black,title={Before you move on}]
What does ruisseler mean? (\textquotedblleft{}To stream, to run down (of water)\textquotedblright{}.) Say it once more.
\end{tcolorbox}
